\documentclass[10pt]{article}

\usepackage[utf8]{inputenc}

\usepackage[T1]{fontenc}

\usepackage{amsmath}

\usepackage{amsfonts}

\usepackage{amssymb}

\usepackage[version=4]{mhchem}

\usepackage{stmaryrd}

\usepackage{graphicx}

\usepackage[export]{adjustbox}

\graphicspath{ {./images/} }



\begin{document}

вом движении, для п. с. за ударной волной, при переменной температуре обтекаемой поверхности.



В П. с. т. трехмерных течений развиты методы решения задач и рассмотрены случаи, приводящие к упрощению уравнений. Уравнения п. с. на скользящем цилиндрич. крыле бесконечного размаха аналогичны уравнениям двумерного п. с. В приближенной постановке решены задачи п. с. на вращающейся цилиндрич. лопасти пропеллера и на вращающемся цилиндре в косо набегающем потоке, а также задача п. с. вблизи линии пересечения двух плоскостей.



Перечисленные исследования п. с. в аэрогидродинамике относятся к первому приближению П. с. т. Высшие приближения позволяют проводить исследования взаимодействия п. с. с внешним потоком, а также расчеты при умеренных числах $R$ (см. [13]).



Устойчивость П. с. позволяет определить границы применимости теории. Имеются исследования (см. [14]) на основе метода малых возмущений с периодическими и локальными начальвыми возмущениями. В случае плоских течений анализ 3-мерных возмущений в линейном приближении сводится на основании теоремы Сквайра к 2-мерному с измекенным значением $v$. Нелинейный анализ при потере устойqивости обнаруживает появление продольных вихрей.



Физич. обобщения задач П. с. т. (см. [15] - [17]) связаны с изучением многофазных течений, с использованием реальных уравнений состояния и коэффициентов переноса (усложнение уравнений), с рассмотрением неравновесных течений с диф̆фуузией (расширение системы уравнений и приобретение ею параболико-гиперболич. типа), с учетом абляции обтекаемой поверхности (усложнение граничных условий и необходимость расчета теплопроводности в теле), с учетом переноса излучения (интегро-дифференциальные уравнения).



Дальнейшее развитие П. с. т. в аәрогидродинамике получила при исследовании течений, не удовлетворяющих предположениям Прандтля, с решениями в виде многослойных асимптотич. разложений. Источ ником сложности структуры решения являются дополнительные малые параметры в граничных условиях (напр., из-за малого радиуса кривизны обтекаемого контура), существования особых точек, линий или поверхностей в первом приближении П. с. т., а также возможная бифуркация решения.



К этому классу относятся течения около точек отрыва и присоединения п. с. к обтекаемому контуру и около точек падения на п. с. ударных волн. Характерные решения (см. [18] - [20]) имеют трехслойную структуру. Внешнее решение определяет возмущенное п. с. потенциальное течение и описывается уравнениями в возмущениях. Средний слой толщины порядка $X \varepsilon$ описывается уравнениями невязких завихренных течений с $p_{y}=0$, и в него поступает газ из основной части предичествующего п. с.



Уравнения 3-го, самого тонкого, пристеночного слоя выводятся в предположении, что его длина и толщина имеют соответственно порядки $\mathrm{X \varepsilon}^{3 / 4}, \mathrm{X \varepsilon}^{5 / 4}$. Введение в стационарном случае переменных



$x=\xi \varepsilon^{-3 / 4}, y=\eta \varepsilon^{-5 / 4}, \quad U=u \varepsilon^{-1 / 4}, V=v \varepsilon^{-3 / 4}, P=p \varepsilon^{-1 / 2}$ приводит при $\varepsilon \rightarrow 0$ снова к уравнениям (3), в к-рых следует $u$ заменить на $U$, а $p$ - на $P$. Эта структура решения задачи характерна для широкого класса течений с малыми величинами возмущений.



Исследованы многие задачи динамики течений при малых $\varepsilon$, в к-рых на коротких расстояниях давление во внетнем сверхзвуковом $(M>1, M=w / c$, где $w-$ скорость газа, с - скорость звука) потоке меняется сильно. $\mathrm{K}$ ним относятся задачи расчета обтекания контуров с большой локальной кривизной п присоеди- нения потока к поверхности тела. В этих случаях при трехслойной схеме решения в среднем слое $p_{y} \neq 0$.



Изучен класс задач, в к-рых возмущенное решение занимает конечную область. Эти решения реализуются на режимах умеренного и сильного взаимодействия п. с. с внешним гиперзвуковым $(M \rightarrow \infty)$ потоком и при сверхзвуковом обтекании тел конечной длины, через поверхность к-рых производится интевсивный вдув $\left(v_{0}>0\right)$ газа. В әтих запачах величина давления на всей внешней границе п. с. определяется из решения полной задачи. Возмущение от задней кромки тела распространяется вверх по потоку, что определяется неединственностью решения в окрестности передней кромки тела.



\includegraphics[max width=\textwidth, center]{2023_07_08_19e8afa9c86f86fad85ag-1(1)}

№ 3, с. $575-86 ;$ [2] B а з о в B., Асимптотические разложес англ., М., 1968; [3] В а с и л в е в а А. Б., Б у т у з о в B. Ф., Асимптотические разложения решений сингулярно возму-



\includegraphics[max width=\textwidth, center]{2023_07_08_19e8afa9c86f86fad85ag-1}

н и к Л. А., "Успехи матем. наук", 1957, т. 12, в. 5, с. 3-122;

1960, т. 15, в. 3, с. 3-80; [5] К о у Л Д ж., Методы возмушений в прикладной математике, пер. с англ., М., 1972; [6] О л е йн и к О. А., «Матем. сб.», 1952, т. $31, \mathcal{N}_{2} 1$, с. 104-17; [7] И л в и в А.' М., ЛІ е л и к о в а E. $\Phi$., там же, 1975, т. 96, № 4, Н. B., Теоретическая гидромеханика, 4 изд., ч. 2 , ї., 1963 ; [9] Јі о й и я н с $\mathrm{k}$ и й ग. Г., Ламинарный пограничный слой, М., 1962; [10] II Л и х ти н г $\Gamma_{. \ddot{ }}$, Теория пограничного слоя, пер. с нем., М., 1974; [11] О ле ин и к О. А., „Успехи матем.



\includegraphics[max width=\textwidth, center]{2023_07_08_19e8afa9c86f86fad85ag-1(2)}

"IІрикл. матем. и механ.», 1942, т. 6, в. 6, с. 449-86; [13] В а нс англ., М., 1967; [14] Б е тч о в Р., К р и м и н а Л е В., Вопросы гидродинамической устойчивости, пер. с англ., М., 1971; [15] С о у С., Гидродинамика многофазных систем, пер. с апгл., М., 1971; [16] Д о р р е н с У. Х Гиперзвуковые течентя вязкого газа, пер. с англ., М., 1966; [17] К ә й с В. М., Конвек-

тивный тепло- и массообмен, пер. с англ., М., 1972; [18] Н е йтивный тепло- и массообмен, пер. с англ., М., 1972; [18] н е и$N_{2} 4$, c. $53-57$; [19] e r o ж e, "Tp. ЦAГИ", 1974, в. 1529; [20] S t e w a r s o n K., "Advances Appl. Mech.", 1974, v. 14, р. ПОГРАНИЧНЫИ СЛОИ - область больших значевий градиента функций, в частности в гидродинамике, область төчевия вязкой жидкости (газа) с малой по сравнению с продольными размерами поперечной толщиной, образующаяся у поверхности обтөкаемого твердого тела или на границе раздела двух потоков жидкостей с различными скоростями, температурами или химич. составами. П. с. характеризуется резким изменением в поперечном направлении скорости (д ин а м и ч е с к и й П. с.) или температуры (те п л 0в о й, и л и т е мп е р а т у р ны й, П. с.), или же концентраций отдельных химич. компонентов (д и ффу у 3 и он ны й, или концент р а ци онны й, П. с.). Понятие П. с. и сам термин были введены Л. Прандтлем (L. Prandtl, 1904) в связи с решением краевой задачи для нелинейных уравнений с частными производными гидродинамики вязких жидкостей. Нужды развития авиации определили разработку теории П. с. в аэрогидродинамике. В сер. 20 в. стала развиваться математическая пограничного слоя теория, а также приложения в теории теплопередачи, дифффузии, процессов в полупроводниках и др. Ю. Д. Шжыллевский.



ПОГРЕШНОСТЬ - разность $x$ - $a$, где $a$ - данное число, к-рое рассматривается как приближенное значение нек-рой величины, точное значение к-рой равно $x$. Разность $x-a$ наз. также а б с о л ю т н о й П. Отношение $x-a$ к $a$ наз. о т н о с и т е л ь н ой П. числа $a$. Для характеристики П. обычно пользуются указанием ее границ. Число $\Delta(a)$ такое, что



\[

|x-a| \leqslant \Delta(a)

\]



наз. г р а в и цей а б с олютн о й П. Число $\delta(a)$ такое, что



\[

\left|\frac{x-a}{a}\right| \leqslant \delta(a)

\]



наз. $\boldsymbol{r} \mathbf{a}$ ци й относительно й П. Гра-





\end{document}